\documentclass[11pt]{article}
\usepackage{amsmath, amssymb}
\usepackage{geometry}
\usepackage{array}
\usepackage{booktabs}
\geometry{margin=0.7in}
\setlength{\parindent}{0pt}
\setlength{\parskip}{0.35em}
\renewcommand{\arraystretch}{1.15}

\newcommand{\cheatsheettitle}[2]{%
\begin{center}
{\LARGE #1}\\[0.4em]
{\large #2}
\end{center}
}


\begin{document}
\cheatsheettitle{Algebra I Exponents and Radicals Cheatsheet}{Module 7}

\section*{Exponent Rules}
For $a\ne0$:
\[
a^m\cdot a^n=a^{m+n}
\]
\[
\frac{a^m}{a^n}=a^{m-n}
\]
\[
(a^m)^n=a^{mn}
\]
\[
(ab)^n=a^nb^n,\qquad \left(\frac{a}{b}\right)^n=\frac{a^n}{b^n}\quad (b\ne0)
\]
\[
a^0=1,\qquad a^{-n}=\frac{1}{a^n}
\]

\section*{Scientific Notation}
\[
a\times 10^n
\]
where
\[
1\le |a|<10
\]
Large numbers have positive powers of $10$. Small decimals have negative powers of $10$.

\section*{Square Roots}
\[
\sqrt{a^2}=|a|
\]
\[
\sqrt{ab}=\sqrt a\sqrt b\qquad (a\ge0,\ b\ge0)
\]
\[
\sqrt{\frac{a}{b}}=\frac{\sqrt a}{\sqrt b}\qquad (a\ge0,\ b>0)
\]

\section*{Simplifying Radicals}
Pull out perfect-square factors.
\[
\sqrt{72}=\sqrt{36\cdot2}=6\sqrt2
\]
\[
3\sqrt5+7\sqrt5=10\sqrt5
\]
Only like radicals combine.

\section*{Rational and Irrational Numbers}
\begin{itemize}
  \item Rational numbers can be written as $\frac{a}{b}$ with integers $a,b$ and $b\ne0$.
  \item Terminating and repeating decimals are rational.
  \item Non-repeating, non-terminating decimals are irrational.
  \item Square roots of non-perfect squares are irrational.
\end{itemize}

\section*{Quick Checks}
\begin{itemize}
  \item Negative exponents move factors across the fraction bar.
  \item Exponents distribute over multiplication, not addition.
  \item $\sqrt{a+b}\ne \sqrt a+\sqrt b$ in general.
  \item Simplified radicals have no perfect-square factors inside.
\end{itemize}

\end{document}
